\chapter{Durchführung}
Zur Kalibrierung der Energie werden die Energiespektren der Nuklide \textsuperscript{22}Na, \textsuperscript{57}Co und \textsuperscript{137}Cs aufgenommen, deren $\gamma$-Energien bekannt sind. Die Breite der Photolinien wird verwendet, um die energieabhängige Energieauflösung zu bestimmen.

Zusätzlich wird das Zeitspektrum mit der \textsuperscript{22}Na-Quelle kalibriert, indem verschiedene Zeitverzögerungen in der Delay-Einheit ein- und ausgeschaltet werden. Danach wird die Verzögerung so eingestellt, dass der Koinzidenzpeak etwa in der Mitte des ADC-Spektrums liegt.

Mit dieser Einstellung wird das Spektrum von \textsuperscript{60}Co aufgenommen. Dabei werden Photolinien, Comptonkanten und Rückstreukanten identifiziert.

Anschließend werden die echten und zufälligen Koinzidenzen mit den \textsuperscript{60}Co-Daten anhand der im vorherigen Abschnitt beschriebenen Methoden ausgewertet. Nun wird noch das 2D-Specktrum bei kurzer Koinzidenzzeit genauer betrachtet. (vgl.   \cite{Vorbereitungshilfe}, S.202)
